\section{习题}
\label{sec:ch1-exercises}
% ============================================================================

\begin{exercise}
\label{ex:ch1-sl2z-1}
验证 $S^2 = -I$，$(ST)^3 = -I$。
由此说明 $\PSL_2(\Z) = \SL_2(\Z)/\{\pm I\}$ 的表现为
$\langle \bar{S}, \bar{T} \mid \bar{S}^2 = (\bar{S}\bar{T})^3 = 1 \rangle$。
\end{exercise}

\begin{exercise}
\label{ex:ch1-fund-domain}
画出 $\GammaO(2)$ 的基本域。
（提示：$[\SL_2(\Z) : \GammaO(2)] = 3$，基本域是 $\SL_2(\Z)$ 基本域的
$3$ 个拷贝的并。）
\end{exercise}

\begin{exercise}
\label{ex:ch1-cusps-gamma0}
证明 $\GammaO(N)$ 在 $\mathbb{P}^1(\Q)$ 上的作用不再传递（当 $N > 1$）。
确定 $\GammaO(N)$ 的尖点等价类个数（用 $N$ 的因子表达）。
\end{exercise}

\begin{exercise}
\label{ex:ch1-weight-k}
证明：若 $k$ 为奇数，则 $\Mk_k(\SL_2(\Z)) = 0$。
（提示：用 $-I \in \SL_2(\Z)$ 的变换律。）
\end{exercise}

\begin{exercise}
\label{ex:ch1-surjection}
证明约化同态 $\SL_2(\Z) \to \SL_2(\Z/N\Z)$ 是满射。
（提示：先证 $N = p$ 素数的情形，利用 Hensel 引理或中国剩余定理。）
\end{exercise}

\begin{exercise}
\label{ex:ch1-isogeny-degree}
设 $\Lambda = \Z\tau + \Z$。证明：$\C/\Lambda$ 到自身的度为 $n$ 的同源
恰好对应于满足 $\alpha\Lambda \subset \Lambda$、$[\Lambda : \alpha\Lambda] = n$
的 $\alpha \in \C$。当 $\Lambda$ 没有复乘（CM）时，说明这样的 $\alpha$ 必为整数。
\end{exercise}
